% CONCLUSIONES Y TRABAJOS FUTUROS

\section{Conclusiones y trabajos futuros}

El desarrollo presentado permitió alcanzar, en un estado funcional inicial, la construcción de un módulo de monitoreo de calidad del aire apoyado en inferencia difusa y datos abiertos. La revisión sistemática del estado del arte permitió caracterizar variables, enfoques difusos y arquitecturas operativas recientes, y aportó la base empírica necesaria para justificar el problema de investigación y derivar criterios explícitos de diseño. Con ello, la propuesta quedó sustentada sobre criterios trazables y coherentes con la evidencia reciente, no sobre decisiones arbitrarias o desvinculadas del campo.

A partir de esa base, se definió e implementó el prototipo \textit{AQRisk}, organizado en cinco capas funcionales: consolidación normativa, variables auxiliares, inferencia difusa principal, capa contextual basada en reglas y alertamiento. La arquitectura implementada integra cálculo del AQI con base en \textit{EPA/AQS}, una base principal de 54 reglas difusas, una base contextual crisp de 9 reglas, control de cobertura de datos, ejecución por CLI y Docker, API HTTP e interfaz web para control, visualización, trazabilidad, explicabilidad y evaluación. Este resultado muestra que la arquitectura funcional propuesta pudo traducirse en un artefacto ejecutable y auditable, con una separación explícita entre la capa normativa, la inferencia difusa principal y el ajuste contextual posterior.

La validación realizada en esta etapa confirma la operabilidad del flujo completo del sistema y la coherencia general de la salida frente a escenarios controlados y consultas reales sobre datos abiertos. La evidencia reunida ya no se limita a pruebas de humo: incluye pruebas automatizadas básicas, corridas reales sobre estaciones con comportamiento útil, evidencia visual por modal y comparación con corridas históricas. El prototipo permite inspeccionar contaminantes soportados, reglas activadas, variables auxiliares, ajuste contextual y diferencias entre episodios reales y controlados. En consecuencia, existe verificación funcional suficiente para sostener la utilidad técnica del artefacto como módulo de apoyo a la evaluación del riesgo y a la generación de alertas, aunque todavía no una evaluación comparativa extensa entre múltiples estaciones, ventanas temporales y condiciones urbanas diversas.

Entre las principales limitaciones del trabajo actual se encuentran la dependencia de disponibilidad y cobertura efectiva en fuentes abiertas vivas, la validación todavía acotada en escenarios reales y la ausencia de una capa completa de gestión de usuarios, roles y administración. Las corridas con datos abiertos también mostraron que no toda estación con sensores publicados ofrece una ventana útil para el modelo, lo que condiciona la comparabilidad entre ciudades y obliga a mantener criterios explícitos de selección de casos. Estas restricciones no invalidan el artefacto, pero sí delimitan el alcance de los resultados presentados: la contribución central del TFM es la construcción y justificación de un módulo explicable, trazable y reproducible para evaluación actual del riesgo, no una plataforma productiva plenamente desplegada a escala urbana.
Entre las principales limitaciones del trabajo actual se encuentran la dependencia de disponibilidad y cobertura efectiva en fuentes abiertas vivas, la validación todavía acotada en escenarios reales y la ausencia de una capa completa de gestión de usuarios, roles y administración. Las corridas con datos abiertos también mostraron que no toda estación con sensores publicados ofrece una ventana útil para el modelo, lo que condiciona la comparabilidad entre ciudades y obliga a mantener criterios explícitos de selección de casos. Estas restricciones no invalidan el artefacto, pero sí delimitan el alcance de los resultados presentados: la contribución central de la propuesta es la construcción y justificación de un módulo explicable, trazable y reproducible para evaluación actual del riesgo, no una plataforma productiva plenamente desplegada a escala urbana.

Como trabajo futuro inmediato, resulta pertinente ampliar la validación comparativa del módulo con más estaciones, más horizontes temporales y casos de alerta de mayor diversidad, además de endurecer la cobertura de pruebas API y E2E. También resulta razonable reforzar la persistencia de corridas con una base de datos más robusta que el histórico local actual, de forma que la comparación temporal deje de depender de almacenamiento básico en el prototipo. En una fase posterior, la implementación podría evolucionar hacia funciones de autenticación, administración y exposición pública controlada, así como hacia una comparación experimental entre la capa contextual crisp actual y una eventual variante contextual difusa.
